\section*{Chapter \ref{ch:dv:trees-and-finite-difference-pricers-in-practice} --- Trees and Finite-Difference Pricers in Practice}

\begin{solution}{exo:dv:trees-and-finite-difference-pricers-in-practice:1}
The tree's terminal nodes sit either on the strike or straddle it depending on the parity of $n$: with the strike on a node the kink is
represented one way, between nodes the other, and the discretisation error of the payoff changes sign. The error is of order $1/n$ with a
parity-dependent constant.
\end{solution}

\begin{solution}{exo:dv:trees-and-finite-difference-pricers-in-practice:2}
$\log_2(0.30/0.091)=1.72$; the next halving should give about $0.091/2^{1.72}=0.028$ cent (the grid gives 0.029).
\end{solution}

\begin{solution}{exo:dv:trees-and-finite-difference-pricers-in-practice:3}
A cash dividend shifts every node by the same amount, not the same proportion, so the up-then-down and down-then-up paths no longer meet and the
tree stops recombining (its size grows exponentially). The grid is fixed in the spot: it applies $V(t_d^-,S)=V(t_d^+,S-D)$ by interpolation and
carries on.
\end{solution}

\begin{solution}{exo:dv:trees-and-finite-difference-pricers-in-practice:4}
With $V_h=V+ch^2$: $V\approx V_{1600}+(V_{1600}-V_{800})/3=6.090331+0.000068/3=6.0903535$.
\end{solution}

\begin{solution}{exo:dv:trees-and-finite-difference-pricers-in-practice:5}
The back substitution must reach the exercise region first, so that each value there is set by the constraint before the continuation values
next to it are computed from it; for a put that region is at the low end, so the substitution runs upwards and the elimination downwards from the
top. Eliminating the other way would compute continuation values from unprojected neighbours.
\end{solution}

\begin{solution}{exo:dv:trees-and-finite-difference-pricers-in-practice:6}
The price's sensitivity to the barrier level is of order one (the derivative of the value with respect to $H$), so an error of half a cell, $O(h)$,
in the level is an $O(h)$ error in the price, which dominates the scheme's $O(h^2)$.
\end{solution}

\begin{solution}{exo:dv:trees-and-finite-difference-pricers-in-practice:7}
Errors of 0.80, 0.35 and 0.16 cent at 100, 200 and 400 nodes: first order (the error halves when the grid doubles), against 0.30, 0.091 and 0.029 with
Brennan--Schwartz. Projecting after an unconstrained Crank--Nicolson solve is inconsistent with the scheme.
\end{solution}

\begin{solution}{exo:dv:trees-and-finite-difference-pricers-in-practice:8}
Repeatability is not accuracy: the tree is deterministic, so it returns the same number every run. At 2\,000 steps its error is still of order 0.04 to
0.08 cent and alternates with the parity of the step count; only a convergence study against a better method certifies the digits.
\end{solution}

\begin{solution}{pb:dv:trees-and-finite-difference-pricers-in-practice:1}
\textbf{1.} 6.0864 and 6.0975; the value is 6.0904.
\textbf{2.} $-0.40$ and $+0.72$ cent.
\textbf{3.} 1\,500 (and 1\,501).
\textbf{4.} $1500\times1501/2=1\,125\,750$.
\textbf{5.} The average reduces the zigzag but keeps the first-order error, whose constant is not known in advance; it has no error estimate.
\textbf{6.} 1.04, 0.30, 0.091 and 0.029 cent.
\textbf{7.} 200 nodes and steps: 40\,600 node updates.
\textbf{8.} Its error at 200 rises to 0.125 cent, over the tolerance; and its \omterm{def:dv:greeks-and-the-hedging-pnl:greeks}{Greeks} near the strike oscillate.
\textbf{9.} 0.11 and 0.054 cent.
\textbf{10.} $V_{200}+(V_{200}-V_{100})/3$, within 0.021 cent of the value.
\textbf{11.} 6.835 European, 7.407 American.
\textbf{12.} The barrier falls between nodes: errors of 0.16 at 200 nodes against 0.0002 with the barrier on a node.
\textbf{13.} Gamma: 0.217 at the strike instead of 0.141 one week from expiry, without Rannacher start-up.
\textbf{14.} A grid node update (a tridiagonal solve) costs a few times a tree node update (a weighted average); the ratio depends on the language and the
implementation.
\textbf{15.} It lowers every method's time by a similar factor, so it does not change which method wins; it moves the whole study to where a million node
updates cost milliseconds.
\textbf{16.} Neither: both are about half a cent off, on opposite sides; use a method whose error is known and small, or at least extrapolate.
\textbf{17.} Price at two grid sizes and report the difference as the error estimate; check against closed forms where they exist; keep the strike and
barriers on nodes; log the grid used.
\textbf{18.} When simplicity matters more than cost (teaching, a quick check), and when exercise rules are easiest to write node by node.
\textbf{19.} Tree: 1\,500 steps, 1\,125\,750 node updates; grid: 200 nodes and steps, 40\,600 node updates; the grid needs 28 times less work.
\textbf{20.} The accurate one has a known, controlled error at the size used; the converging one only promises to get there eventually.
\end{solution}

\begin{solution}{iq:dv:trees-and-finite-difference-pricers-in-practice:1}
Crank--Nicolson damps high-frequency error components only weakly; a kinked payoff injects them at the strike, and near expiry there are too few
steps for them to decay. Gamma, a second derivative, amplifies them. Fix: a few fully implicit (Rannacher) half steps at the start, and \omterm{def:dv:trees-and-finite-difference-pricers-in-practice:smoothing}{payoff
smoothing}.
\iqlookfor{weak damping, the kink, and the two remedies.}
\end{solution}

\begin{solution}{iq:dv:trees-and-finite-difference-pricers-in-practice:2}
Crank--Nicolson in $\ln S$ with the strike on a node, Rannacher start-up, and the \omterm{def:dv:trees-and-finite-difference-pricers-in-practice:bs}{Brennan--Schwartz algorithm} for the constrained tridiagonal system (one
sweep per step), checked by two grid sizes; not projection after an unconstrained solve, which falls to first order.
\iqlookfor{Brennan--Schwartz and the accuracy details.}
\end{solution}

\begin{solution}{iq:dv:trees-and-finite-difference-pricers-in-practice:3}
Binomial: first order, with an oscillation between odd and even step counts. Trinomial: first order, smoother. Crank--Nicolson with smoothing, start-up and
the strike on a node: second order. For a tenth of a cent the grid needs about 30 times less work than the binomial tree in the chapter's case.
\iqlookfor{the orders and the oscillation.}
\end{solution}

\begin{solution}{iq:dv:trees-and-finite-difference-pricers-in-practice:4}
Dividend: a \omterm{def:dv:trees-and-finite-difference-pricers-in-practice:jump}{jump condition} at the ex-date, $V(t_d^-,S)=V(t_d^+,S-D)$, by interpolation, then a Rannacher restart. Barrier: a grid whose boundary (or a node) is
exactly at the barrier, with the rebate as the boundary value; discrete monitoring as a projection at the dates.
\iqlookfor{\omterm{def:dv:trees-and-finite-difference-pricers-in-practice:jump}{jump condition} and alignment.}
\end{solution}

\begin{solution}{iq:dv:trees-and-finite-difference-pricers-in-practice:5}
Convergence studies at several grid sizes with the observed order compared with theory; agreement with closed forms and with independent methods (a tree, a
Monte Carlo); limit tests (zero volatility, far strikes); \omterm{def:dv:greeks-and-the-hedging-pnl:greeks}{Greeks}' smoothness; regression tests pinned to certified numbers; and the same checks for the
production (C++) and research (Python) versions.
\iqlookfor{order of convergence, independent references, regression.}
\end{solution}

\begin{solution}{iq:dv:trees-and-finite-difference-pricers-in-practice:6}
Candidates: the dividend's \omterm{def:dv:trees-and-finite-difference-pricers-in-practice:jump}{jump condition} applied at the wrong step or without a Rannacher restart (oscillations near the boundary); interpolation across the
ex-date that smears the kink of the exercise value; the Brennan--Schwartz direction not reversed for a call (exercise region at high spots); the ex-date not
on the time grid. Check each against a fine grid and a tree with the dividend.
\iqlookfor{the dividend handling and the elimination direction.}
\end{solution}
